Approximate nearest-neighbour (ANN) search reduces query time by avoiding
exhaustive, full-precision distance evaluation. Graph methods such as CAGRA
restrict visits through navigation, while compressed IVF combines selective
list access with inexpensive scoring \cite{cagra,johnson_gpu,ivfrabitq}.
The tested embeddings have 768--3,072 coordinates and up to 17.78 million
vectors; Stella's raw FP32 array alone requires 67.81\,GiB
(\cref{sec:7-7}). Compact codes address capacity, but table construction,
lookups, selection and launch overhead determine their query-time benefit.

JHQ separates screening from refinement: compact primary codes score
candidates, and residuals refine a shortlist \cite{jhq}. This exposes two
decisions: how to execute the broad screen and how many survivors to refine.
The trained index can be reused while varying execution or budget, allowing
their effects to be tested separately. This thesis evaluates these decisions
through Recall@10, host-visible throughput/latency, resident memory and
construction cost, including configurations where a competitor or the
unmodified design is faster.
